\section*{Capítulo \ref{ch:b3:sensory-systems} --- Sistemas sensoriales}

\begin{solution}{exo:b3:sensory-systems:1}
Por el cable, no por el mensaje: cada nervio sensorial es una \omterm{def:b3:sensory-systems:principles}{línea
marcada} que termina en su propia región del encéfalo, y todo lo que
llega por el nervio óptico se lee como luz ---por eso la presión sobre
el globo ocular produce un destello, y un golpe en el oído, un zumbido.
\end{solution}

\begin{solution}{exo:b3:sensory-systems:2}
En la oscuridad, el GMP cíclico mantiene abiertos los canales
catiónicos y una ``corriente de oscuridad'' entrante constante mantiene
la célula despolarizada y liberando glutamato. La luz activa la
\omterm{def:b3:sensory-systems:retina}{rodopsina}, la transducina y la fosfodiesterasa, que destruye el GMPc;
los canales se cierran, la corriente entrante se detiene y la célula se
hiperpolariza y libera menos ---la luz se señaliza mediante una
disminución.
\end{solution}

\begin{solution}{exo:b3:sensory-systems:3}
$I = I_{0}10^{L/10}$: \qty{20}{dB}, \qty{1e-10}{W/m^2}; \qty{60}{dB},
\qty{1e-6}{W/m^2}; \qty{110}{dB}, \qty{0.1}{W/m^2}. Del más fuerte al
más suave: $10^{9}$.
\end{solution}

\begin{solution}{exo:b3:sensory-systems:4}
El dulce, el umami y el amargo a través de receptores acoplados a
proteínas~G (uno del dulce, uno del umami y unos veinticinco del
amargo); el salado a través del canal de sodio ENaC; el ácido a través
de un canal de protones. La mayor parte del ``sabor'' es el olor de
moléculas volátiles que alcanzan la nariz desde la boca; con la nariz
taponada solo quedan los cinco \omterm{def:b3:sensory-systems:chemical}{gustos} y la comida parece sosa.
\end{solution}

\begin{solution}{exo:b3:sensory-systems:5}
Fechner: $S = (1/0.02)\ln(10\,000/100) = 50\ln 100 \approx 230$
escalones. Contando: $1.02^{n} = 100$ da $n = \ln 100/\ln 1.02 = 233$.
\end{solution}

\begin{solution}{exo:b3:sensory-systems:6}
$P = 1 - e^{-a}\sum_{k=0}^{5}a^{k}/k!$: $a = 2$: $0.017$; $a = 6$:
$0.55$; $a = 12$: $0.98$. La mitad de los destellos se ven en
$a \approx 5.7$; con una absorción del \qty{10}{\%}, el destello debe
entregar unos $57$ fotones en la córnea.
\end{solution}

\begin{solution}{exo:b3:sensory-systems:7}
$55/3\times 1.3 \approx 24$; en \omterm{def:b3:sensory-systems:cochlea}{decibelios} de presión, $20\log_{10}24
\approx \qty{28}{dB}$. Unos huesecillos fusionados no transmiten la
vibración, de modo que el sonido llega a la \omterm{def:b3:sensory-systems:cochlea}{cóclea} solo por conducción
a través del cráneo, sin la ganancia del oído medio: una pérdida de
transmisión de unos \qtyrange{30}{50}{dB}.
\end{solution}

\begin{solution}{exo:b3:sensory-systems:8}
El canal de la \omterm{def:b3:sensory-systems:cochlea}{célula ciliada} lo abre directamente el \omterm{def:b3:sensory-systems:cochlea}{enlace apical} al
tirar de él ---sin mensajero, sin enzima---, de modo que responde en
microsegundos; el fotorreceptor amplifica un solo fotón a través de dos
etapas enzimáticas, que tardan decenas de milisegundos. La audición
gana la precisión temporal necesaria para localizar sonidos por los
retrasos interaurales y para seguir la fase; la visión gana la
sensibilidad para contar fotones aislados, y lo paga en rapidez.
\end{solution}

\begin{solution}{exo:b3:sensory-systems:9}
$H(0.05) = -0.05\ln 0.05 - 0.95\ln 0.95 = 0.150 + 0.049 = 0.199$. Ser
humano: $\ln\binom{400}{20} \approx 400\times 0.199 = 80$, unos
$10^{34}$ patrones. Ratón: $\ln\binom{1100}{55} \approx 1100\times 0.199
= 219$, unos $10^{95}$.
\end{solution}

\begin{solution}{exo:b3:sensory-systems:10}
Media de $500\times 0.1/40 = 1.25$ sucesos espontáneos por ventana. $P(k
\le 5) = e^{-1.25}(1 + 1.25 + 0.78 + 0.33 + 0.10 + 0.03) = 0.998$, de
modo que $P(k \ge 6) \approx 2\times 10^{-3}$. Con un umbral de uno, en
la mayoría de las ventanas ocurriría un suceso espontáneo
($1 - e^{-1.25} = 0.71$) y la oscuridad estaría llena de destellos; con
seis, las falsas alarmas llegan una vez cada quinientas ventanas. El
umbral lo dicta el ruido de los bastones, no su sensibilidad.
\end{solution}

\begin{solution}{exo:b3:sensory-systems:11}
Retraso máximo, $0.2/340 = \qty{590}{\micro s}$ (sonido procedente de un
lado). Cerca de la línea media $\Delta t \approx (d/c)\sin\theta$, de
modo que $\qty{10}{\micro
s}$ corresponde a $\sin\theta = 0.017$, alrededor de
\qty{1}{\degree}. Por encima de \qty{4}{kHz} el período
(\qty{250}{\micro s}) es más corto de lo que las neuronas pueden
enganchar, y un retraso de varios cientos de microsegundos abarca más
de un ciclo, lo que vuelve ambigua la fase; el encéfalo emplea entonces
la diferencia de nivel entre los oídos y el código de lugar.
\end{solution}

\begin{solution}{exo:b3:sensory-systems:12}
Lejos del borde, $r = I(1 - 2w)$: $0.6$ en el lado oscuro y $1.2$ en el
claro. En el último receptor oscuro ($I = 1$, vecinos $1$ y $2$): $r
= 1 - 0.2\times 3 = 0.4$; en el primero claro: $r = 2 - 0.6 = 1.4$. El
escalón queda exagerado en un defecto y un exceso ---las bandas de
Mach. La \omterm{def:b3:sensory-systems:retina}{retina} gana contraste y bordes, y un intervalo dinámico menor
que transmitir; pierde fidelidad a los niveles absolutos y produce
ilusiones de brillo.
\end{solution}

\begin{solution}{pb:b3:sensory-systems:1}
\textbf{1.} $0.01/3.6\times 10^{-19} = 2.8\times 10^{16}$ fotones por
segundo.
\textbf{2.} Área de la pupila, $\pi(3.5\times 10^{-3})^{2} = 3.8\times 10^{-5}$
m$^{2}$. A \qty{1}{km}: fracción $3.8\times 10^{-5}/1.26\times 10^{7} =
3\times 10^{-12}$, $8.5\times 10^{4}$ fotones por segundo. A
\qty{10}{km}: $850$ por segundo.
\textbf{3.} En \qty{0.1}{s} con el \qty{10}{\%} absorbido: $850$ a
\qty{1}{km} y $8.5$ a \qty{10}{km} ---ambos por encima de seis; la vela
se ve a diez kilómetros (la mayor parte de las veces).
\textbf{4.} Seis absorbidos por ventana exigen $600$ fotones por segundo
dentro de la pupila: $4\pi d^{2} = 3.8\times 10^{-5}\times 2.8\times 10^{16}/600
= 1.8\times 10^{9}$ m$^{2}$, $d \approx \qty{12}{km}$. Allí, $P(k \ge
6 \mid a = 6) = 0.55$: la vela se ve en alrededor de la mitad de las
ventanas.
\textbf{5.} $e^{-6} = 0.0025$. En el umbral el recuento fluctúa de una
ventana a otra ---$6 \pm 2.4$---, de modo que la fuente parece ir y
venir: el centelleo de una estrella tenue (al que la atmósfera añade el
suyo propio).
\textbf{6.} La fluctuación del fondo, $\sqrt{10^{4}} = 100$ fotones por
ventana, ahoga los $8.5$ de la vela: la detección exige que la señal
supere el ruido del fondo, no el umbral del ojo adaptado a la
oscuridad.
\textbf{7.} $\qty{1}{pA}/200 = \qty{5}{fA}$ por canal; $5\times
10^{-15}/1.6\times 10^{-19} = 3\times 10^{4}$ iones por segundo.
\textbf{8.} Una treintava parte; unos treinta fotones simultáneos
cierran todos los canales y saturan el \omterm{def:b3:sensory-systems:retina}{bastón}.
\textbf{9.} $10^{-12}\times 0.2 = 2\times 10^{-13}$ C, $1.3\times 10^{6}$
cationes: un fotón gobierna más de un millón de cargas.
\textbf{10.} El \omterm{def:b3:sensory-systems:retina}{bastón} se satura en milisegundos, con todos los canales
cerrados, y su \omterm{def:b3:sensory-systems:retina}{rodopsina} se blanquea más deprisa de lo que se regenera;
los \omterm{def:b3:sensory-systems:retina}{conos}, de ganancia menor y apagado más rápido, siguen respondiendo
y desplazan su intervalo a lo largo de seis décadas de luz.
\textbf{11.} Una ganancia alta exige una integración larga (cada etapa
lleva tiempo y la respuesta se prolonga para acumularse), de modo que
el \omterm{def:b3:sensory-systems:retina}{bastón} es lento y sensible; la amplificación menor y el apagado más
rápido del \omterm{def:b3:sensory-systems:retina}{cono} dan una respuesta en \qty{20}{ms}, bastante rápida para
el movimiento, a costa de necesitar un centenar de fotones.
\textbf{12.} $10^{8}/40 = 2.5\times 10^{6}$ isomerizaciones espontáneas
por segundo en toda la \omterm{def:b3:sensory-systems:retina}{retina} ---pero solo $1.25$ por ventana en
cualquier grupo de $500$ bastones, muy por debajo del umbral de seis,
de modo que el ruido se rechaza grupo a grupo; el umbral y la
agrupación existen precisamente para esto.
\textbf{13.} \qty{0.1}{W/m^2}; sobre \qty{5.5e-5}{m^2}, \qty{5.5e-6}{W}.
\textbf{14.} $5.5\times 10^{-6}\times 7200 = \qty{0.04}{J}$ ---unas
cuatrocientas gotas de lluvia a lo largo de dos horas.
\textbf{15.} $10^{-12}\times 5.5\times 10^{-5} = 5.5\times 10^{-17}$ W;
en \qty{10}{ms}, $5.5\times 10^{-19}$ J ---aproximadamente la energía de
uno o dos fotones visibles.
\textbf{16.} $0.3/5000 = 6\times 10^{-5}$ rad, $0.003^{\circ}$; la
desviación equivale a unos tres átomos de hidrógeno.
\textbf{17.} $10^{(110 - 85)/10} = 316$: dos horas a \qty{110}{dB}
entregan la energía de $630$ horas a \qty{85}{dB} ---casi un mes de
jornadas laborales.
\textbf{18.} Unos $120$ escalones apenas perceptibles, uno por \omterm{def:b3:sensory-systems:cochlea}{decibelio}.
\textbf{19.} $2^{400} = 10^{400\log_{10}2} = 10^{120}$.
\textbf{20.} $\ln\binom{400}{20} \approx 400\,H(0.05) \approx 80$, menos
una corrección de Stirling de alrededor de $2$: unos $10^{34}$ patrones.
\textbf{21.} Cuéntense los receptores presentes en un patrón y no en el
otro: $5 +
5 = 10$ de $40$ (una distancia de Hamming); los patrones que se solapan
en tres cuartas partes activan casi los mismos \omterm{def:b3:sensory-systems:chemical}{glomérulos} y se leen
como casi el mismo olor.
\textbf{22.} $10^{12}/10^{34} = 10^{-22}$ de los patrones. El límite no
es el código: los receptores son ruidosos, muchos están ajustados de
forma parecida y sus actividades están correlacionadas, las moléculas
reales producen solo un subconjunto pequeño de patrones, y la
discriminación y la memoria de patrones del encéfalo son finitas.
\textbf{23.} $\ln\binom{1100}{20} \approx 1100\,H(0.018) = 1100\times
0.091 \approx 100$, unos $10^{43}$ ---unas $10^{9}$ veces la cifra
humana de patrones de veinte receptores (y muchos más si el ratón
emplea más receptores por olor).
\textbf{24.} El olfato general apenas cambia ---el código es redundante
y otros receptores captan todos los odorantes---, pero un odorante que
dependa de ese receptor a concentración baja se vuelve indetectable o
cambia de carácter: una anosmia específica, como la de las muchas
personas que no huelen la androstenona.
\textbf{25.} La vela se ve hasta unos \qty{12}{km}; el sonido umbral
entrega $5.5\times 10^{-19}$ J al tímpano en \qty{10}{ms}, uno o dos
fotones; una nariz humana puede formar unos $10^{34}$ patrones de
veinte receptores.
\end{solution}
